\chapter{Open problems}\label{ch:open}

\noindent\emph{Editorial note (September 2026).} This list was compiled on 28 August 2026,
before Chapters \ref{ch:XXVII} and \ref{ch:XXVIII} were written and before the September
line-by-line review of the corpus. Several items are now settled and are kept here for the
record only: condition (C4) and every spectral-gap route to it are resolved, in the
negative, in Chapter \ref{ch:XXVIII} (the orbit relation is hyperfinite for every
admissible monoid); the non-abelian intermediate structure is resolved in Chapter
\ref{ch:XXVII} (the ordered Frobenius cocycle realizes $\mathrm{Prob}(H\backslash G)$, with
lift-dependence and the rigidity criterion); and the sieve-family items should be read
against the revised Chapter \ref{ch:XXIV} (two obstructed families, not four). Where an
item below conflicts with those chapters, the chapters are authoritative.\medskip

\begin{enumerate}
\item \textbf{A non-abelian model between the two walls.} Does there exist an Ore
      semigroup that is non-abelian but whose range projections are nested rather than
      partitioning? Abelian semigroups cannot produce non-abelian essential values
      (Chapter~\ref{ch:VII}); free non-abelian semigroups empty the high-temperature region
      (Chapter~\ref{ch:V}); the Hurwitz quaternions satisfy (C1)--(C3)
      (Chapter~\ref{ch:X}). \emph{[Settled: condition (C4) fails for every admissible
      monoid, nested or not --- Chapter~\ref{ch:XXVIII}. The residual question is whether
      any index structure outside the (C1)--(C3) framework, necessarily of infinite norm
      at some place, can carry both arithmetic thermodynamics and non-amenability.
      Chapter~\ref{ch:XXXIV} now closes the scaling route universally (the abelian
      scaling bottleneck: scalar dynamics on a right LCM semigroup cannot produce a
      nonabelian residual symmetry) and opens the angular route: the angular Hurwitz
      cocycle $a\mapsto a/\mathrm{Nrd}(a)^{1/2}\in SU(2)$ has Kakutani abscissa exactly
      $2$, and the expected verdict $\mathrm{KMS}_\beta\cong\mathrm{Prob}(SU(2))$ for
      $\beta>2$ is reduced to one named transport problem.]}
\item \textbf{Condition (C4).} Is the measured groupoid of the Hurwitz system non-amenable?
      The Ramanujan route is closed by the nested/partition dichotomy
      (Chapter~\ref{ch:XI}); the question reduces to a spectral gap for the single
      representation $\pi_0|_{G^1}$ of the norm-one subgroup, together with essential
      freeness of the $G^1$-action (Chapter~\ref{ch:XII}). \emph{[Settled, in the
      negative: the orbit relation is hyperfinite, the factor injective, and attaining the
      Ramanujan bound certifies temperedness, the signature of amenability ---
      Chapter~\ref{ch:XXVIII}.]}
\item \textbf{Cartan uniqueness without norm-separation.} Is $C(Y_S)$ intrinsically
      characterised inside $(\cA_{K,S},\sigma)$ when the relation lattice is nontrivial?
      The Nica--Toeplitz relations block the obvious counterexample
      (Chapter~\ref{ch:XIV}); the spectral route cannot decide it
      (Chapter~\ref{ch:XV}); mixed isometries exist iff a single measure-preserving
      transformation is not ergodic (Chapter~\ref{ch:XV}, Prop.~3.2).
\item \textbf{Which field is recovered.} The equivariant $C^*$-system sees
      $(S,\Ns|_S,G_S)$; whether this triple determines $K$ is an arithmetic question
      (Chapter~\ref{ch:XIII}). Over function fields the answer is known to stop at the
      isogeny class plus a torsor coordinate (Chapters~\ref{ch:XIX}--\ref{ch:XX}).
\item \textbf{The exact value of $W_1(\varphi_+,\varphi_-)$} (Chapter~\ref{ch:IV}).
      \emph{[Partially settled: Chapter~\ref{ch:XXIX} reduces the problem to a sequence of
      second-order cone programs with certified two-sided error, and replaces the squeeze
      lower bound by the strictly sharper quadratic chaos bound
      $\Pi_\infty(w_0^2+4\sum_jw_j^2t_j^2/(1-t_j)^2)^{1/2}$. Remaining: (a) does the odd
      Walsh-chaos hierarchy exhaust $W_1$, or do higher-degree terms strictly increase the
      supremum; (b) the matching upper bound --- now advanced in Chapter~\ref{ch:XXXIII}: the adaptive
      coupling gives the closed form $w_0\Pi_\infty+\sum_kw_{j_k}\tau_{j_k}\prod_{l<k}m_{j_l}$,
      strictly below $w_0$ and pinching against the squeeze lower bound in the deep
      freeze; the residual factor is now fully diagnosed in Chapter~\ref{ch:XXXVII}: the single-step
      Kantorovich polytope is exactly the scan bound (fractional transport gains
      nothing), diagonal displacements carry the closed-form cone discount
      $d_L=w(1+1/\sqrt2)<2w$, and the one remaining instrument is the discrete
      $\infty$-Laplacian geodesic problem for $d_L$; (c) the optimal profile is NOT parity-measurable
      (numerical verdict, Chapter~\ref{ch:XXIX}, Remark 4.3).]} Also open: a transport
      quantity built from the tail $\sigma$-algebra that would detect the transition
      (Chapter~\ref{ch:XVI}).
\item \textbf{$K_*(\partial\cA^{\mathrm{aff}})$} of the affine boundary
      (Chapter~\ref{ch:II}). \emph{[Partially settled: Chapter~\ref{ch:XXX} computes both
      groups at a principal single prime --- the norm survives as the exact order $q-1$ of
      $[1]$, the trace of the uniformizer enters the intermediate exterior layers, and the
      invariant hears the unit orbit of the uniformizer, hence is not canonical in the
      prime. Remaining after Chapter~\ref{ch:XXXI}'s ray-class repair (which proves the class-group
      order $h$ is the exact index of definable scalings and promotes the order of $[1]$
      to $N\pp^{\,h}-1$): (a) separating $q$ from $h$: generically resolved --- the intermediate layers separate,
      e.g.\ $\Q(\sqrt{-5})/\pp_3$ vs $\Q(i)/(3)$ share $\mathrm{ord}[1]=8$ but have
      $T_1$ of orders $2$ vs $4$ --- with only layer-coincident pairs left to the dual
      coaction; (b) the Galois descent from the Hilbert class
      field --- Chapter~\ref{ch:XXXII} now proves the full class group emerges as the
      group of descent obstructions $\Cl_K\hookrightarrow H^1(G,\mathcal O_H^\times)$,
      with the ladder $\mathrm{ord}[1]=q^h-1$ (bottom) and $q^{|\Cl_K|}-1$ (top), and
      reduces the equivariant K-theory to the inert local model crossed by $\Z/h$; the
      integral computation of that reduction is now carried out for the first inert case
      in Chapter~\ref{ch:XXXVIII} (exact torsion by Smith forms, the unit class split by
      $\tfrac{1\pm U_\sigma}2$ into a pair of order-$8$ classes, and the entanglement
      law: the class group doubles only the non-induced strata); remaining there are the
      twisted character bookkeeping and the degree-four extension data; (c) the
      several-prime case, where the definable scalings are the kernel of
      $\Z^S\to\Cl_K$; plus the extension data for $n\ge3$ and the $q$-primary
      bookkeeping from Chapter~\ref{ch:XXX}.]} Also open:
      $H^1_{\mathrm{meas}}(\mathcal R_S,\pi_\beta;\T)$ itself rather
      than its kernel (Chapter~\ref{ch:III}). \emph{[Settled in the correct sense by
      Chapter~\ref{ch:XXXIX}: the group is non-smooth and its coboundary action is
      turbulent, so no classification by countable structures exists; the complete
      smooth residue is the Mackey-range spectrum; the arithmetic embeds as
      $\widehat{G_S}/\Xi_\beta$, the unique equicontinuous thread, while the ambient
      group is orbit-equivalence amnesiac. Remaining: the flagged turbulence lemma and
      noncompact coefficients.]}
\item \textbf{The regime of the twin primes} in the free system, equivalent to comparing a
      Brun-type constant with $1$ (Chapter~\ref{ch:V}).
      \emph{[Settled, unconditionally: a rigorous segmented-sieve computation
      (Chapter~\ref{ch:XXXV}, \texttt{code/twin-regime/}) gives $\sum_{p\le1.1\times10^{10}}1/p=1.00196>1$
      over the twin primes --- a partial sum of a positive series, hence an unconditional
      lower bound --- so $L_{S}(1)>1$ and the free system is in Regime (S),
      $\beta_c^{\mathrm{free}}>1$. The exact value of $\beta_c^{\mathrm{free}}$ remains
      conditional on the tail.]}
\item \textbf{A knot family that is linking-rich and volume-sparse}, which the arithmetic
      topology dictionary requires and which is not known (Chapter~\ref{ch:VIII}).
\item \textbf{Whether a refinement of the obstruction data recovers the essential-value
      subgroup $E_\beta$ up to conjugacy} in the Galois-twisted groupoid
      (Chapter~\ref{ch:VII}).
\item \textbf{Closed phase and no congruence obstruction simultaneously.} A closed
      low-temperature phase needs density zero with two prime conditions; an unobstructed
      high-temperature phase is easiest for trace conditions, which have positive density.
      The Chen primes are the only known candidate satisfying both
      (Chapters~\ref{ch:XXIV}--\ref{ch:XXV}).
\end{enumerate}

\medskip\noindent
Two items are open at the level of the project rather than the mathematics and are
recorded in the repository's status file: a repeatable priority search against the
literature on KMS states of groupoid $C^*$-algebras with infinite sparse $S$, and an
independent verification of the bibliography.
